% =====================================================================
% TASK 3 — SENTIMENT CLASSIFICATION
% Paste this section into the main report (before Task 5, after the PMI
% section). Requires: \usepackage{booktabs}, graphicx (already loaded).
% Confusion-matrix figures are in results/task3_outputs/.
% =====================================================================

\section{Sentiment Classification}
\label{sec:classification}

This section builds classical classifiers for the sentiment task and
analyses their behaviour. We train a Multinomial Naive Bayes (NB)
classifier and a feed-forward neural network (a multilayer perceptron,
MLP), each with two feature representations, and we study the effect of
class imbalance, of dropping the neutral class, and of the typical
errors. The aim, following the EDA, is not to maximise accuracy but to
understand \emph{why} each modelling choice changes the results.

\subsection{Setup}

The dataset was split into 80\% training and 20\% test data with a
stratified split (\texttt{random\_state=42}), which preserves the
59/28/12 class proportions in both parts. Stratification is important
here because the negative class is only 12\% of the data: a non-stratified
split could place too few negatives in the test set and make the
metrics unstable. All classical models use the \texttt{text\_stopwords}
preprocessing variant. Features were built only from the training set
(\texttt{fit} on train, \texttt{transform} on test) to avoid information
leakage. We report accuracy, per-class precision/recall, macro-averaged
F1, and confusion matrices; as argued in the EDA, macro-F1 is the
primary metric because accuracy is dominated by the majority class.

\subsection{Naive Bayes: Bag-of-Words vs.\ TF-IDF}

We first compare two representations for Naive Bayes: bag-of-words
counts (BoW) and TF-IDF. The training BoW matrix has 8733 features and
is 99.9\% sparse, consistent with the short headlines and large
vocabulary found in the EDA. Table~\ref{tab:nb-rep} reports the results.

\begin{table}[t]
\centering
\caption{Naive Bayes with two representations and with a uniform prior.}
\label{tab:nb-rep}
\begin{tabular}{lccc}
\toprule
Configuration & Accuracy & Macro-F1 & Neg.\ recall \\
\midrule
NB + BoW                & 0.72 & 0.66 & 0.53 \\
NB + TF-IDF             & 0.68 & 0.45 & 0.06 \\
NB + BoW (uniform prior) & 0.68 & 0.64 & 0.63 \\
\bottomrule
\end{tabular}
\end{table}

Counter to the common assumption that TF-IDF is a stronger
representation, it \emph{degraded} Naive Bayes sharply: macro-F1 fell
from 0.66 to 0.45 and the recall on the negative class collapsed to
0.06. The mechanism is specific to NB. NB scores a class as the product
of a class prior and per-word likelihoods. TF-IDF down-weighting
flattens the word-likelihood signal, so the high prior of the neutral
class (0.59) dominates and the model defaults to predicting neutral.
This is a clear example of the EDA warning that accuracy can be
misleading: accuracy moved only slightly (0.72 to 0.68), while macro-F1
revealed that the model had effectively stopped predicting the minority
classes.

\subsection{Handling Class Imbalance}

Because both representations showed the model collapsing toward the
neutral majority, we tested an explicit imbalance correction. Naive
Bayes has no \texttt{class\_weight} parameter, but setting
\texttt{fit\_prior=False} replaces the data-estimated priors with a
uniform prior, removing the majority-class advantage. With BoW
(Table~\ref{tab:nb-rep}, last row), negative recall rose from 0.53 to
0.63 and the three classes became markedly more balanced
(all F1-scores between 0.57 and 0.77), at the cost of a small drop in
accuracy (0.72 to 0.68). Macro-F1 itself barely changed (0.66 to 0.64),
because the gain on the rare classes offset the loss on neutral. The
value of this correction is therefore not a higher headline score but a
model that no longer ignores the minority class.

\subsection{Feed-forward Neural Network}

We then trained an MLP (one hidden layer of 100 ReLU units, Adam,
\texttt{max\_iter=200}, \texttt{random\_state=42}; roughly 876{,}000
parameters) on the same representations. Table~\ref{tab:grid} places all
four model$\times$representation combinations side by side.

\begin{table}[t]
\centering
\caption{Macro-F1 for the full classifier $\times$ representation grid.}
\label{tab:grid}
\begin{tabular}{lcc}
\toprule
 & BoW & TF-IDF \\
\midrule
Naive Bayes      & 0.66 & 0.45 \\
Neural network   & 0.67 & 0.65 \\
\bottomrule
\end{tabular}
\end{table}

Two findings stand out. First, with BoW the neural network only
\emph{matched} Naive Bayes (macro-F1 0.67 vs.\ 0.66), despite its much
larger capacity and its ability to model word interactions. This
indicates that for these short headlines word \emph{presence} is nearly
sufficient and word \emph{interactions} add little: knowing whether
words such as \textit{loss}, \textit{profit}, or \textit{growth} are
present already captures most of the signal. Second, and more
interestingly, TF-IDF devastated Naive Bayes (0.45) but left the neural
network almost unaffected (0.65; negative recall 0.54 vs.\ NB's 0.06).
The reason is that the neural network has no class prior to be
overwhelmed --- it learns its weights directly by gradient descent. This
shows that representation effects are \emph{not absolute}: they interact
with the classifier. A study that ran only NB with TF-IDF would have
wrongly concluded that ``TF-IDF is worse''; running the full grid
reveals that TF-IDF is worse \emph{for Naive Bayes specifically}.

\subsection{Error Analysis}

We examined the misclassified test examples of the best multiclass model
(NN + BoW, 267 errors out of 970). The confusion matrices
(Figure~\ref{fig:cm-nb}) show that errors are overwhelmingly
confusions with the neutral class rather than polarity flips: positives
and negatives are rarely confused with each other, but both are often
predicted as neutral. The difficulty is therefore separating weak
sentiment from neutral reporting, not telling positive from negative.

Reading the most frequent error type --- true positives predicted as
neutral --- reveals a consistent pattern: these headlines contain almost
no overtly positive \emph{words}. We group the failures into four
categories:

\begin{itemize}
  \item \textbf{Sentiment in events/numbers:} e.g.\ ``\textit{subscriptions
  increase Cargotec's share capital\dots}'' or ``\textit{awarded a
  significant order for 292 cranes}''. The signal is the financial event
  (a capital increase, a won contract), not a sentiment word.
  \item \textbf{Implicit/connotative:} e.g.\ ``\textit{moving to new
  headquarters signifies the beginning of a new era}''. Positive by
  implication, with no positive keywords.
  \item \textbf{Negation / polarity reversal:} e.g.\ ``\textit{Stora
  Enso\dots was \underline{acquitted} of charges of participating in a
  price-fixing \underline{conspiracy}}''. The sentence is full of
  negative words; \textit{acquitted} reverses them, but bag-of-words
  ignores word order.
  \item \textbf{Domain knowledge required:} e.g.\ being listed among
  ``\textit{sustainable / pioneer companies}'' is an honour that requires
  world knowledge to read as positive.
\end{itemize}

These categories are a direct consequence of the data properties found
in the EDA --- short, domain-specific headlines where sentiment is often
implicit --- and of the bag-of-words assumption that discards word order.
They also motivate the transformer experiments in
Section~\ref{sec:transformers}: each failure mode (word order for
negation, world knowledge for events, context for connotation) is
something a pretrained contextual model is designed to capture.

\begin{figure}[t]
\centering
\includegraphics[width=0.49\linewidth]{../results/task3_outputs/Naive_Bayes___BoW.png}
\includegraphics[width=0.49\linewidth]{../results/task3_outputs/Neural_Net___BoW.png}
\caption{Confusion matrices (multiclass, BoW): Naive Bayes (left) and
neural network (right). Both confuse positives and negatives with the
neutral majority rather than with each other.}
\Description{Two 3x3 confusion matrices showing most off-diagonal mass in the neutral column.}
\label{fig:cm-nb}
\end{figure}

\subsection{Binary vs.\ Multiclass}

Finally, we repeated the task in a binary setting by removing the neutral
class (1967 headlines, positive 69\% / negative 31\%). Removing neutral
also reduces the imbalance: the negative class rises from 12\% of the
data to 31\%. Table~\ref{tab:binary} compares both settings.

\begin{table}[t]
\centering
\caption{Macro-F1 in the multiclass and binary settings.}
\label{tab:binary}
\begin{tabular}{lcc}
\toprule
Model & Multiclass & Binary \\
\midrule
NB + BoW & 0.66 & 0.81 \\
NN + BoW & 0.67 & 0.83 \\
\bottomrule
\end{tabular}
\end{table}

Both models improve substantially (macro-F1 about $+0.15$), and negative
recall recovers strongly (NB 0.53 to 0.70, NN 0.57 to 0.74). The
interpretation matters more than the numbers: every error category above
was a confusion \emph{with neutral}, so deleting the neutral class
removes precisely the dominant failure mode. What remains --- separating
positive from negative --- is the part bag-of-words already handled well.
The $\sim$0.15 gain therefore quantifies how much of the multiclass
difficulty was due to the neutral class alone: the binary setting removes
the hard part, not the easy part. The relative behaviour of the two
models is unchanged (the neural network again only marginally beats Naive
Bayes), confirming that for this dataset word presence is nearly
sufficient and a simple, interpretable, instantly-trained Naive Bayes is
a strong and arguably preferable baseline.
